\clearpage
\VerdierPhysicalPageAnchor{0239}%
\begingroup\small\setlength{\parskip}{0pt}\setlength{\emergencystretch}{1.5em}
\setlength{\abovedisplayskip}{4pt plus 2pt minus 2pt}%
\setlength{\belowdisplayskip}{4pt plus 2pt minus 2pt}%

\VerdierSemanticAnchor{ch4d-4-4-2-b-cont}
\noindent\textit{On a alors :}
\[
\mathrm I_2^{p,q}=H_{\mathcal B}^pF^{q,\bullet}(Y^\bullet).
\]

\VerdierSemanticAnchor{ch4d-4-4-2-c}
\noindent\textit{c) Les termes $\mathrm I_\infty^{p,q}$ sont les gradués
associés aux objets $F^{p+q}(Y)$ convenablement filtrés.}

\VerdierSemanticAnchor{ch4d-4-4-2-proof}
Démontrons a). Soit $\bigl(Y(p,q),\delta(p,q,r)\bigr)$ le premier objet
spectral à valeurs dans $K(\mathcal A)$ associé au complexe $Y$. Il résulte
immédiatement de la définition de cet objet qu'on a (4.1.2.3) :
\[
Y(-p-1,-p)=Y^p[-p].
\]
Par suite (chap.~II, 4.3.3.1) :
\[
\mathrm I_1^{p,q}=F^{p+q}\bigl(Y^p[-p]\bigr)=F^q(Y^p).
\]
La description de la différentielle
$d_1\colon\mathrm I_1^{p,q}\longrightarrow\mathrm I_1^{p+1,q}$ se déduit
alors immédiatement de (4.1.2.6). L'assertion b) résulte immédiatement de
a). La dernière assertion résulte de (chap.~II, 4.3.2).

\VerdierSemanticAnchor{ch4d-4-4-3}
\medskip\noindent\textbf{Remarque 4.4.3.}\quad
Désignons par
$\bigl(\mathrm I_r^{p,q}(Y),d_r,r\geq1\bigr)$ la première suite spectrale
du foncteur $F$ relative au complexe $Y$. Le foncteur :
\[
Y\longmapsto\bigl(\mathrm I_r^{p,q}(Y),d_r,r\geq1\bigr)
\]
est un foncteur de la catégorie $\operatorname{comp}(\mathcal A)$ dans la
catégorie des suites spectrales à valeurs dans $\mathcal B$. Ce foncteur ne
se factorise pas en général par la catégorie $K(\mathcal A)$. En revanche,
le foncteur :

\VerdierSemanticAnchor{ch4d-4-4-3-r-ge-2}
\[
Y\longmapsto\bigl(\mathrm I_r^{p,q}(Y),d_r,r\geq2\bigr)
\]
se factorise, d'après la proposition (4.4.2), b), par la catégorie des
complexes à homotopie près.

\VerdierSemanticAnchor{ch4d-4-4-4}
\medskip\noindent\textbf{4.4.4.}\quad
Soient maintenant $\mathcal A$ une catégorie abélienne,
$F\colon D(\mathcal A)\longrightarrow\mathcal B$ un foncteur cohomologique
à valeurs dans une catégorie abélienne $\mathcal B$, $F^p$, $p\in\mathbb Z$,
les foncteurs déduits de $F$ en composant avec les translations de
$D(\mathcal A)$ ($F^p(X)=F(X[p])$). Soient, de plus, $Y$ un objet de
$\operatorname{comp}(\mathcal A)$ et
$\bigl(Y(p,q),\delta(p,q,r)\bigr)$ le premier objet spectral, à valeurs dans
$D(\mathcal A)$, associé au complexe $Y$. Le foncteur $F$ transforme cet
objet spectral en un objet spectral à valeurs dans $\mathcal B$. On en
déduit une suite spectrale à valeurs dans $\mathcal B$ qu'on appelle encore
la \emph{première suite spectrale du foncteur $F$, relative au complexe
$Y$}. Cette suite spectrale est décrite essentiellement par la proposition
(4.4.2) (il suffit dans la proposition (4.4.2) de remplacer la catégorie
$K(\mathcal A)$ par la catégorie $D(\mathcal A)$). La première suite
spectrale du foncteur $F$, relativement au complexe $Y$, dépend
fonctoriellement de $F$ et de l'argument $Y$ quand celui-ci varie dans
$\operatorname{comp}(\mathcal A)$. Nous utiliserons toujours, pour cette
suite spectrale, l'indexation introduite en (chap.~II, 4.3.3.1).
\endgroup
